\documentclass{article}
\usepackage{fullpage,fourier,amsmath,amssymb}
\usepackage{listings,color,url,hyperref}
\usepackage{epigraph, graphicx}
\usepackage[x11names]{xcolor}

\title{Assignment 1 \\ Left, Right and Center}
\author{Prof. Darrell Long \\
CSE 13S -- Spring 2021}
\date{Due: April 11$^\text{th}$ at 11:59\,pm}

\usepackage{fancyhdr}
\pagestyle{fancy}
\fancyhf{}

\fancypagestyle{plain}{%
  \fancyhf{}
  \renewcommand{\headrulewidth}{0pt}
  \renewcommand{\footrulewidth}{0pt}
  \lfoot{\textcopyright{} 2021 Darrell Long}
  \rfoot{\thepage}
}

\pagestyle{plain}

\definecolor{codegreen}{rgb}{0,0.5,0}
\definecolor{codegray}{rgb}{0.5,0.5,0.5}
\definecolor{codepurple}{rgb}{0.58,0,0.82}

\lstloadlanguages{C,make,python,fortran}

\lstdefinestyle{c99}{
    morekeywords={bool, uint8_t, uint16_t, uint32_t, uint64_t, int8_t, int16_t, int32_t, int64_t},
    commentstyle=\color{codegreen},
    keywordstyle=\color{magenta},
    numberstyle=\tiny\color{codegray},
    identifierstyle=\color{blue},
    stringstyle=\color{codepurple},
    basicstyle=\ttfamily,
    breakatwhitespace=false,
    breaklines=true,
    captionpos=b,
    keepspaces=true,
    numbers=left,
    numbersep=5pt,
    showspaces=false,
    showstringspaces=false,
    showtabs=false,
    tabsize=4
}

\lstset{language=C, style=c99}

\newcommand\asgn[1]{asgn#1}

\begin{document}\maketitle

%%%%%%%%%%%%%%%%%%%%%%%%%%%%%%%%%%%%%%%%%%%%%%%%%%%%%%%%

\section{Introduction}
\epigraphwidth=0.5\textwidth
\epigraph{\emph{The gambling known as business looks with austere disfavor
upon the business known as gambling.}}{---Ambrose Bierce}

\noindent We are going to implement a simple game called \emph{Left, Right, and Center}.
A subset of philosophers will represent players of the game.
It requires no skill, no real decision-making, and a philosopher that is out of the
game can suddenly come back in and often win.

Some of you may have religious or moral prohibitions against gambling. This
program is not gambling, since (i) you neither win nor lose, and (ii) only
fictitious people lose or win any money. But, if you do have any qualms, let
us know and we will give you an alternative assignment involving vampires.

\section{Playing the Game}
\epigraphwidth=0.75\textwidth
\epigraph{\emph{No sympathy for the devil; keep that in mind. Buy
the ticket, take the ride\ldots and if it occasionally gets a little
heavier than what you had in mind, well\ldots maybe chalk it off
to forced conscious expansion: Tune in, freak out, get beaten.}}{---Hunter
S. Thompson, \emph{Fear \& Loathing in Las Vegas}}\noindent

Some number of $k$ players, $1 < k \le 14$, sit around a table.
Each player has in her hand \$3. There are three dice, and each die
has $6$ faces and is labeled: $3 \times \textbf{\textbullet}$, $1
\times \textbf{L}$, $1 \times \textbf{R}$ or $1 \times \textbf{C}$.
As a result, we know that there is a $50$\% chance of rolling
\textbf{\textbullet}, and $16.6\overline{6}$\% chance of rolling
each of \textbf{L}, \textbf{R}, or \textbf{C}.

\begin{enumerate}
\item
Beginning with player $1$, roll the dice:
\begin{enumerate}
\item
If the player has \$3 or more then she rolls three dice;
if she has \$2 then she rolls two dice;
if she has only \$1 then she rolls one die;
if she has no money then she must pass.

\item For each die:
\begin{enumerate}
\item If the player rolls \textbf{L} then she gives \$1 to the player on her \emph{left}.

\item If the player rolls \textbf{R} then she gives \$1 to the player on her \emph{right}.

\item If the player rolls \textbf{C} then she puts \$1 in the pot in the
    \emph{center}.

\item If the player rolls \textbf{\textbullet} then she ignores it.
\end{enumerate}

\end{enumerate}
\item Move to the next player in sequence: to the
right. The players are numbered. There is \texttt{you}. Then there
is the left player which is (\texttt{you} $-$ 1) \texttt{mod} the number
of players and there is the right player which is (\texttt{you} $+$
1) \texttt{mod} the number of players.
Be careful: What does \texttt{-2 \% 10} mean? Consequently, you may find this code useful:
\begin{codelisting}{}
//
// Returns the position of the player to the left.
//
// pos:     The position of the current player.
// players: The number of players in the game.
//
static inline uint8_t left(uint8_t pos, uint8_t players) {
  return ((pos + players - 1) % players);
}

//
// Returns the position of the player to the right.
//
// pos:     The position of the current player.
// players: The number of players in the game.
//
static inline uint8_t right(uint8_t pos, uint8_t players) {
  return ((pos + 1) % players);
}
\end{codelisting}

\item Repeat until only one player has any money remaining (who then wins the pot).
\end{enumerate}

%%%%%%%%%%%%%%%%%%%%%%%%%%%%%%%%%%%%%%%%%%%%%%%%%%%%%%%%

\section{Your Task}
\epigraphwidth=0.6\textwidth
\epigraph{
\emph{Any one who considers arithmetical methods of producing random digits
is, of course, in a state of sin.}}{---John von Neumann, 1951}\noindent

\begin{itemize}

\item You must have one source file: \texttt{lrc.c}. Do not name
your source file anything else. \textcolor{red}{You will lose points}.


\item For grading purposes the numbering of the faces matters, and
so they must be defined as follows (do not change it):

\begin{codelisting}{}
typedef enum faciem { PASS, LEFT, RIGHT, CENTER } faces;
faces die[] = { LEFT, RIGHT, CENTER, PASS, PASS, PASS };
\end{codelisting}

\textcolor{red}{This means you may not use \texttt{int} as a substitute.}

\item You must give your players names, and for grading purposes the names
must correspond to these (do not change them):

\begin{codelisting}{}
char *philosophers[] = { "Immanuel Kant",
                         "Martin Heidegger",
                         "David Hume",
                         "Georg Wilhelm Friedrich Hegel",
                         "Arthur Schopenhauer",
                         "Ludwig Wittgenstein",
                         "Karl Wilhelm Friedrich Schlegel",
                         "Friedrich Nietzsche",
                         "Socrates",
                         "John Stuart Mill",
                         "Plato",
                         "Aristotle",
                         "Thomas Hobbes",
                         "Rene Descartes" };
\end{codelisting}

This array of names is defined in \texttt{philos.h}, which can be found in
the course resources repository.

\item Typing \texttt{make} must build your program and
\texttt{./lrc} must run your program. Since you have not learned
about \texttt{Makefiles} yet, here is one that you can use for now.

\begin{codelisting}{\texttt{Makefile}}
CC = clang
CFLAGS = -Wall -Wextra -Werror -Wpedantic

all: lrc

lrc: lrc.o
    $(CC) -o lrc lrc.o

lrc.o: lrc.c
    $(CC) $(CFLAGS) -c lrc.c

clean:
    rm -f lrc lrc.o

scan-build: clean
    scan-build make
\end{codelisting}

\item You will need to use a random number generator to simulate
rolling a dice. You can do this by calling the function \texttt{random()}
(read the \emph{man page}) to get a random value. You can then use \texttt{mod}
to limit this value to the range of 0--5 inclusive (in Computer Science, we start from 0).
This value is used to determine what was rolled.

\item In order that your program be \emph{reproducible}, you must start from a
known place. This is accomplished by \emph{setting the random seed}
using the function \texttt{srandom()} (again read the \emph{man page}).
Your program will ask for two numbers: the random seed and the
number of players. You should assign these inputs to variables to use
in your program. The random seed completely determines the outcome of your program. If you give it
the same random seed and the number of players you \emph{must} get
the same answer. Here is an example of using \texttt{srandom()} and
\texttt{random()}:
\begin{codelisting}{}
srandom(1);       // Sets the random seed as 1.
int a = random(); // Set a as a pseudorandom number.
\end{codelisting}

\item \emph{Comment your code.} Include block comments to tell the grader what a certain block of code does. Use a line comment to explain something that is not as obvious.
\textcolor{red}{Refer to the coding standards}.

\item All source code should be formatted using \texttt{clang-format}.

\item Your program \emph{must} compile and run on Ubuntu 20.04 with no
  errors. If it does not, your program will receive a 0.

\item Your program must pass \texttt{scan-build} with no errors.
  Document any errors that are false positives in your \texttt{README.md}.

\end{itemize}

In lecture we talked briefly about \emph{random} and \emph{pseudo-random}
numbers. As we know, computers produce pseudo-random numbers,
and in this case it is to your benefit since \emph{reproducibility}
is essential. That means that in reality you program though it
appears to be random is actually \emph{deterministic}. This is why
starting with the same seed produces the same sequence.

%%%%%%%%%%%%%%%%%%%%%%%%%%%%%%%%%%%%%%%%%%%%%%%%%%%%%%%%
\section{Resources}
\epigraph{\emph{Common sense in an uncommon degree is what the
world calls wisdom.}}{---Samuel Taylor Coleridge, \emph{Literary Remains}}\noindent

A working example binary can be found in the course resources repository
on \texttt{git.ucsc.edu}. The header file, \texttt{philos.h}, containing
the names of the philosophers can be found here as well.

For this assignment, running ``\texttt{./lrc}'' will prompt the user to
enter a random seed and the number of players. It will then produce an
output of the game which your program should also be able to produce
given the same seed and number of players. Your program must be able to
do \textit{exactly} the same.

%%%%%%%%%%%%%%%%%%%%%%%%%%%%%%%%%%%%%%%%%%%%%%%%%%%%%%%%
\section{Hints}
\epigraph{\emph{It is no use trying to sum people up. One must follow hints,
not exactly what is said, nor yet entirely what is done.}}{---Virginia Woolf}\noindent

\begin{enumerate}

\item \emph{\textcolor{red}{Start Now!}} You can always finish early.

\item You may find it helpful to draw out and run through a game to get a feel
    for the rules. Doing so may make it easier to visualize the game and design
    your program.
\item The game itself should be an \textit{infinite} loop, where the condition to break
    out of this loop is when there is only one active player remaining.
\item You should think carefully about what quantities that you
must track in order for your program to function. At a \emph{minimum}
you must keep track of the bank balance of each player, the amount
of money in the pot, and the number of players that are \emph{in}.
Be careful: players that were \emph{out} may be brought back in if
money is passed to the \emph{left} or \emph{right}.

\end{enumerate}

%%%%%%%%%%%%%%%%%%%%%%%%%%%%%%%%%%%%%%%%%%%%%%%%%%%%%%%%

\section{Deliverables}
\epigraph{\emph{The design process is about designing and prototyping and
making. When you separate those, I think the final result
suffers.}}{---Jonathan Ive}\noindent

You will need to turn in:

\begin{enumerate}
\item \texttt{lrc.c}: The source file which is your program.
\item \texttt{philos.h}: The header file containing the names of the
  philosophers.
\item \texttt{Makefile}: This is a file that will allow the grader to type
    \texttt{make} to compile your program. Typing \texttt{make} with no
    arguments must build your program.
\begin{itemize}
\item \texttt{CC=clang} must be specified.
\item \texttt{CFLAGS=-Wall -Wextra -Werror -Wpedantic} must be included.
\item \texttt{make clean} must remove all files that are compiler generated.
\item \texttt{make} should build your program, as should \texttt{make all}.
\item Your program executable must be named \texttt{lrc}.
\end{itemize}

\item \texttt{README.md}: This must be in markdown.
This must describe how to use your program and \texttt{Makefile}.

\item \texttt{DESIGN.pdf}: This must be in PDF format. The design document
should describe your design for your program with enough detail
that a sufficiently knowledgeable programmer would be able to
replicate your implementation. This does not mean copying your
entire program in verbatim. You should instead describe how your
program works with supporting pseudo-code. For this program, pay
extra attention to how you describe your logic flow/algorithm in
\textbf{C}. \textcolor{red}{The commit ID containing the first draft of
  your design document must be submitted no later than the Thursday
  before the assignment is due, by 11:59 pm PST. Failing to submit this
  first draft or submitting a insubstantial draft means receiving a 0
for your design document.}

\end{enumerate}

All of these files must be in the directory \texttt{\asgn{1}}.

%%%%%%%%%%%%%%%%%%%%%%%%%%%%%%%%%%%%%%%%%%%%%%%%%%%%%%%%

\section{Submission}
\epigraphwidth=0.5\textwidth
\epigraph{\emph{Better three hours too soon than a minute too
late.}}{---William Shakespeare}

To submit your assignment, refer back to \texttt{\asgn{0}} for the steps on how to submit your assignment through \texttt{git}. Remember: \emph{add, commit,} and \emph{push}!

\textcolor{red}{Your assignment is turned in \emph{only} after you have pushed.
If you forget to push, you have not turned in your assignment and you will get
a \emph{zero}. ``I forgot to push'' is not a valid excuse. It is \emph{highly} recommended to commit and push your changes \emph{often}.}
%\end{enumerate}

%%%%%%%%%%%%%%%%%%%%%%%%%%%%%%%%%%%%%%%%%%%%%%%%%%%%%%%%
\section{Supplemental Readings}
\epigraph{\emph{I was not proud of what I had learned but I never doubted that it was worth knowing.}}{---Hunter S. Thompson, \emph{The Rum Diary}}\noindent

 \begin{itemize}

 	\item \textit{The C Programming Language} by Kernighan \& Ritchie
	\begin{itemize}
		\item Chapter 1 \S 1.6, 1.7
		\item Chapter 2 \S 2.3
		\item Chapter 3 \S 3.1-3.7
	\end{itemize}
 \end{itemize}

%%%%%%%%%%%%%%%%%%%%%%%%%%%%%%%%%%%%%%%%%%%%%%%%%%%%%%%%%%%

\begin{center}
\includegraphics[width=0.5\textwidth]{../../images/monkey-chainsaw.jpg}\\
\emph{datis mandatis ut simia scriptura gladium}
%\emph{البرمجة بلغة C كإعطاء منشار آلي لقرد}
%\emph{La programmation en \textbf{C} c'est comme donner une tron\c{c}onneuse \`{a} un singe.}
%\emph{用\textbf{C}语言编程就像给猴子一把电锯一样。}
%\emph{Is cos\'{u}il le r\'{i}omhchl\'{a}r\'{u} i \textbf{C} s\'{a}bh slabhrach a thabhairt do monca\'{i}.}
\end{center}


\end{document}



